\section*{Wstęp}
\addcontentsline{toc}{section}{Wstęp}

Rozwój technologii analizy obrazu wykorzystujących modele uczenia maszynowego umożliwia coraz szersze zastosowanie systemów informatycznych do oceny ruchu człowieka. Jednym z obszarów zastosowania tych rozwiązań jest sport, a w szczególności analiza techniki wykonywania ćwiczeń na podstawie parametrów biomechanicznych. Prawidłowe wykonanie ćwiczenia ma istotne znaczenie zarówno dla skuteczności treningu, jak i ograniczenia ryzyka kontuzji wynikającego z błędów technicznych.

Martwy ciąg jest jednym z podstawowych wielostawowych ćwiczeń siłowych i należy do konkurencji trójboju siłowego. Ćwiczenie to wymaga skoordynowanej pracy wielu segmentów ciała, w szczególności stawów biodrowych i kolanowych oraz odpowiedniej kontroli położenia tułowia. Choć ruch może wydawać się prosty, w rzeczywistości jest złożony i wymaga odpowiedniej techniki wykonania. Błędy techniczne podczas wykonywania martwego ciągu mogą prowadzić do przeciążeń układu ruchu. Z tego względu analiza techniki wykonania stanowi istotny element treningu.

Tradycyjna ocena techniki wykonywania ćwiczeń siłowych opiera się na obserwacji prowadzonej przez trenera lub inną osobę posiadającą odpowiednie kompetencje, a także na samodzielnej obserwacji wykonywanego ruchu, na przykład w lustrze. Taka ocena zależy między innymi od doświadczenia osoby oceniającej i może być trudna do przeprowadzenia w sposób powtarzalny. Zastosowanie metod komputerowej analizy obrazu pozwala na częściową automatyzację tego procesu. Umożliwia powtarzalne wyznaczanie mierzalnych parametrów opisujących ruch oraz wskaźników pozwalających podsumować technikę wykonania ćwiczenia.

W pracy wykorzystano gotowy model estymacji pozy człowieka, umożliwiający wykrywanie punktów kluczowych sylwetki odpowiadających położeniom wybranych stawów. Na podstawie położenia tych punktów możliwe jest obliczanie wybranych parametrów biomechanicznych, między innymi kątów w stawach oraz zmian położenia wybranych punktów w czasie. Parametry te są następnie wykorzystywane do oceny wybranych elementów techniki oraz spójności kolejnych powtórzeń.

Celem pracy jest zaprojektowanie i zaimplementowanie aplikacji sieciowej umożliwiającej analizę biomechaniki martwego ciągu na podstawie nagrania wideo. System wykorzystuje gotowy model estymacji pozy człowieka do wykrywania punktów kluczowych sylwetki, a następnie oblicza wybrane parametry ruchu i prezentuje wyniki użytkownikowi za pośrednictwem interfejsu aplikacji klienckiej.

Praca ma charakter aplikacyjny. Zakres pracy nie obejmuje opracowania nowego modelu uczenia maszynowego, lecz wykorzystanie istniejącego rozwiązania do stworzenia kompletnego systemu informatycznego integrującego moduł analizy nagrań wideo z interfejsem użytkownika. System umożliwia przesłanie nagrania, przeprowadzenie automatycznej analizy oraz przedstawienie jej wyników w czytelnej formie.